\documentclass[12pt]{article}
\usepackage[utf8]{inputenc}
\usepackage[brazil]{babel}
\usepackage[T1]{fontenc}

\begin{document}

\section{Metas e Atividades Realizadas de Maio}

\subsection{Objetivos propostos para o mês:}
\begin{enumerate}
    \item \textbf{Estruturação de Dados:} Configuração do ambiente Django e criação dos modelos fundamentais para usuários, disciplinas e semestres no banco de dados.
    \item \textbf{Base de Interface:} Inicialização do projeto em React Native com a definição da navegação base e implementação das telas de autenticação.
\end{enumerate}

\subsection{Atividades realizadas no mês:}

\subsubsection{Configuração Base do Repositório}
O repositório oficial no GitHub foi criado e configurado seguindo a arquitetura de monorepositório, estabelecendo as pastas raiz backend e frontend para manter o desenvolvimento das equipes isolado e organizado. Todos os integrantes foram integrados como colaboradores.

\subsubsection{Divisão Técnica da Equipe}
Para garantir o paralelismo entre a inteligência do servidor e a interface visual, a equipe foi segmentada em duas frentes de trabalho:
\begin{itemize}
    \item \textbf{Frontend (3 integrantes):} Responsáveis pela criação da interface fluida, componentes reutilizáveis e fluxo de navegação. Composta pelos alunos Julio, Nicolas Fenato e Laura Brenny
    \item \textbf{Backend (2 integrantes):} Focados na modelagem de dados, segurança da API e lógica de processamento no servidor. Composta pelos alunos Rafael Pinheiro e Eduardo Trautwein.
\end{itemize}
\subsubsection{Fundação do Backend}
O ambiente do servidor foi inicializado localmente com o isolamento de um ambiente virtual venv para evitar conflitos de bibliotecas. Realizou-se a instalação estruturada do framework Django e do Django REST Framework, criando os arquivos bases de gerenciamento do sistema como o manage.py e a pasta de configurações globais. O motor do backend está preparado para suportar a lógica das tabelas e a futura transmissão de dados via API.
\subsubsection{Fundação do Frontend}
O projeto base em React Native foi inicializado dentro da pasta frontend utilizando as ferramentas oficiais de criação do framework. A estrutura de arquivos gerada foi devidamente sincronizada e submetida ao repositório oficial no GitHub, que agora esta devidamente pronto para o desenvolvimento das telas e dos fluxos de navegação mobile.

\subsubsection{Modelagem e Mapeamento de Dados}
A equipe concluiu a definição do Modelo Entidade-Relacionamento, estabelecendo as tabelas fundamentais. Utilizando o sistema de ORM (Object-Relational Mapping) do Django, essas entidades foram mapeadas para o banco de dados, permitindo que as tabelas físicas sejam geridas diretamente via código Python.

\subsubsection{Criação da Navegação Base e Tema Global} 
A estrutura de rotas do frontend foi feita com a implementação do pacote React Navigation, dividindo o fluxo do aplicativo em duas frentes organizacionais. O fluxo de autenticação inicial foi estruturado através do arquivo AppNavigator.js, que atua como um navegador em pilha Stack Navigator responsável por gerenciar as transições entre as telas de login e cadastro. Ademais, a estrutura das telas internas foi feita por meio do TabNavigator.js, utilizando Bottom Tab Navigation, que é responsável por projetar o menu de abas inferiores com os direcionamentos para os módulos de principais (calendário, dashboard e disciplinas). Além disso, para assegurar um padrão estético no desenvolvimento, foi criado um arquivo theme.js, que centraliza as padronizações globais da paleta de cores, tipografia, espaçamento e arredonadamento das bordas.

\subsubsection{Implementação das Telas de Autenticação}
A criação das interfaces de acesso foi iniciada com o estabelecimento dos esqueletos da LoginScreen.js e da CadastroScreen.js na pasta de telas do projeto. Essa estrutura viabilizou a validação imediata do fluxo de rotas em pilha gerenciado pelo AppNavigator.js, evitando erros de importação no sistema.


\section{Metas e Atividades Realizadas de Junho}

\subsection{Observações do mês:}
Devido à grande quantidade de trabalhos e provas no mês de junho, a equipe teve disponibilidade reduzida para dedicação ao projeto. Porém, graças ao adiantamento realizado anteriormente, os módulos de Faltas e Calendário já se encontram parcialmente estruturados no backend, conforme detalhado a seguir.

\subsection{Objetivos propostos para o mês:}
\begin{enumerate}
    \item \textbf{Frequência:} Desenvolvimento dos algoritmos de cálculo de presença e alertas de limite de faltas na interface.
    \item \textbf{Cronograma Dinâmico:} Criação dos modelos de eventos, provas e trabalhos no backend e desenvolvimento da visualização de calendário no mobile.
\end{enumerate}

\subsection{Atividades realizadas no mês:}

\subsubsection{Backend}
Como adiantamento realizado ainda em maio, o backend dos módulos de Faltas e Calendário já está bem estruturado: as entidades falta, avaliacao, tarefa e notificacao foram totalmente mapeadas via ORM do Django, com seus respectivos relacionamentos de chave estrangeira com a disciplina, controle de unicidade para evitar lançamentos duplicados de faltas, cálculo de peso para médias ponderadas e vinculação de notificações a tarefas ou avaliações.

\begin{verbatim}
    
class Falta(models.Model):
    id_falta = models.AutoField(primary_key=True)
    id_disciplina = models.ForeignKey(Disciplina, ...)
    data_aula = models.DateField()
    quantidade = models.IntegerField()

class Avaliacao(models.Model):
    id_avaliacao = models.AutoField(primary_key=True)
    id_disciplina = models.ForeignKey(Disciplina, ...)
    titulo = models.CharField(max_length=150)
    peso = models.DecimalField(max_digits=5, decimal_places=2)
    nota = models.DecimalField(max_digits=5, decimal_places=2, null=True)

class Tarefa(models.Model):
    id_tarefa = models.AutoField(primary_key=True)
    id_disciplina = models.ForeignKey(Disciplina, ...)
    id_avaliacao = models.ForeignKey(Avaliacao, null=True)
    descricao = models.CharField(max_length=255)
    prazo = models.DateField()
    concluida = models.BooleanField()

class Notificacao(models.Model):
    id_notificacao = models.AutoField(primary_key=True)
    id_tarefa = models.ForeignKey(Tarefa, null=True)
    id_avaliacao = models.ForeignKey(Avaliacao, null=True)
    mensagem = models.CharField(max_length=255)
    agendada_para = models.DateTimeField()
\end{verbatim}

\subsubsection{Frontend}
Com o backend à frente do cronograma, a meta para o próximo mês é que a equipe de frontend desenvolva as telas e componentes visuais para acompanhar o que já foi estruturado no servidor, priorizando o controle de frequência e a visualização do calendário acadêmico.

\end{document}
